\section{The connected rank-three triangular locus}

Fix the original frame $F$ satisfying $F.\mathrm{Ok}(C_0)$,
with odd residue characteristic $\ell$. Write $k=O/\mathfrak m_O$,
$G_0=\ker\delta$, $g^*=cg^{-1}c^{-1}$ and $p_g=\Psi(g)$.
Then $p_{g^*}=p_g^{-1}$. For an original rank-three space $S$,
write $R$ for its complete universal ring, $B$ for the closed
weight-zero coefficient ring and $A$ for the $B$-algebra of
matrix entries. The weights in order are $(1,0,-1)$.
Write $\mathfrak t_+$ for the ideal of strictly lower entries,
$\mathfrak t_-$ for that of strictly upper entries, and $I_E$
for the ideal of entries of odd weight. The locus under study is
the source's
\[
  \mathrm{BB}=\{\mathfrak p\in\operatorname{Spec} A:
      \mathfrak m_B A\subseteq\mathfrak p,\quad
      \mathfrak t_+\not\subseteq\mathfrak p\}.
\]
We construct its connecting curves in this spectrum.

% BEGIN G67T1
\paragraph{G67T1. A normalized lower triangular point with a nonzero adjacent entry.}
For a lower unitriangular matrix write
\[
 L(\alpha,\beta,\gamma)=
 \begin{pmatrix}1&0&0\\\alpha&1&0\\\gamma&\beta&1\end{pmatrix}.
\]
Apply RTC11 to $\chi=\Psi^{-1}\delta_k$. This character takes
value $-1$ at $c$, the frame's local separators have $p_s^2\ne1$,
and its $\mathrm{f4}$ condition is precisely the nonvanishing of
the required strict cohomology. We get a locally constant
$1$-cocycle $x$ on the full group, zero at $c$ and nonzero somewhere
on $G_0$. Define $u_g=p_gx(g)$ for $g\in G_0$. The cocycle rule,
$x(c)=0$, and the local coboundary formula of RTC11 give
\[
 u_{gh}=p_hu_g+u_h,\qquad u_{g^*}=p_g^{-1}u_g,\qquad
 u_s=(p_s-1)A_v\quad(s\text{ in the local image at }v)
\]
for scalars $A_v\in k$. Indeed on $G_0$ the coefficient action is
$p_g^{-1}$, so $x(g^{-1})=-p_gx(g)$ and
$x(cg^{-1}c^{-1})=-x(g^{-1})$. The function $u$ is nonzero.

Consider the following matrices on $G_0$ and an element $D$ of
the lower unipotent group:
\[
 M(g)=
 \begin{pmatrix}p_g&0&0\\u_g&1&0\\
 u_g^2/(2p_g)&u_g/p_g&p_g^{-1}\end{pmatrix},
 \qquad
 d=-u_{g_0}/(p_{g_0}+1),\quad D=L(d,-d,-d^2/2).
\]
Both denominators are units, since the mark has $p_{g_0}^2\ne1$.
The rule for $u_{gh}$ gives $M(gh)=M(g)M(h)$, with the bottom-left
identity obtained by expanding $(p_hu_g+u_h)^2/(2p_gp_h)$.
Reversal of the transpose gives $M(g^*)=JM(g)^tJ$ for the
antidiagonal matrix $J$. Since $JD^tJ=D^{-1}$, conjugating every
matrix by $D$ preserves this identity. The new adjacent entries
are $y_g=u_g+d(p_g-1)$ and $y_{g^*}=u_g/p_g+d(p_g^{-1}-1)$.
At $g_0$ they are equal by the formula for $d$, giving the marked
normalization. If $y$ vanished identically, its values at $g_0$
and $g_0^*$ and the formula for $u_{g_0^*}$ would imply
$2d(p_{g_0}-1)=0$. Thus $d=0$, contradicting $u\ne0$.

On the local group at $v$, set $L_v=L(A_v,A_v,A_v^2/2)$. Then
\[
             M(s)=L_v\operatorname{diag}(p_s,1,p_s^{-1})L_v^{-1}.
\]
Thus one lower unitriangular basis
works for that entire local group, including after multiplication
by $D$. Let $a_g$ denote the new bottom-left entry. We have produced
a pair $(y_*,a_*)$ of locally constant lower-entry functions
satisfying all the normalized group and ordinary identities,
with $y_*\ne0$.
% END G67T1

% BEGIN G67T2
\paragraph{G67T2. Lower entries and the ordered ordinary condition.}
In any commutative $k$-algebra $Q$ put
\[
 M_{y,a}(g)=\begin{pmatrix}p_g&0&0\\y_g&1&0\\
              a_g&y_{g^*}&p_g^{-1}\end{pmatrix}.
\]
The source condition $F.\mathrm{Low}(Q,y,a)$ says that $y,a$ are
locally constant and satisfy
\[
 \begin{aligned}
 y_{gh}&=p_h y_g+y_h,\qquad y_{g_0}=y_{g_0^*},\\
 a_{gh}&=p_h a_g+y_{g^*}y_h+p_g^{-1}a_h,\qquad a_{g^*}=a_g .
 \end{aligned}
\]
Its remaining condition at each $v$ is that, for fixed
$\alpha,\beta,\gamma\in Q$ there,
\[
 M_{y,a}(s)=L(\alpha,\beta,\gamma)
       \operatorname{diag}(p_s,1,p_s^{-1})L(\alpha,\beta,\gamma)^{-1}
                    \qquad(s\in\operatorname{im}\operatorname{loc}_{0,v}).
\]
Multiplication gives the three lower entries
$(p_s-1)\alpha$, $(1-p_s^{-1})\beta$, and
$(p_s-p_s^{-1})\gamma-(1-p_s^{-1})\alpha\beta$ at positions
$(1,0),(2,1),(2,0)$. Thus the display is the exact three local
equations. These formulas give a normalized representation with
prescribed diagonal and local characters. In particular the ordered
ordinary product vanishes for every local tuple, since in that fixed
local basis each diagonal entry of the product has a zero factor.
They apply to the pair of G67T1.

Conversely, take a lower triangular normalized representation with
locally constant entries, and the specified residual diagonal and
local characters in $Q$. Polarization under $g\mapsto g^*$ determines
the $(2,1)$ entry from the $(1,0)$ entry and fixes $(2,0)$.
The globally displayed identities follow by multiplication and
marked normalization. At $v$ choose a separator $s_0$ from the
frame. Its three diagonal values $p,1,p^{-1}$ have unit pairwise
differences, since $p^2\ne1$ in $k$. Successive substitution in a
lower triangular matrix gives a lower unitriangular eigenbasis for
this element over $Q$. In that basis let its diagonal entries
be $d_i$. In the ordinary product use $s_0$ in every position except
position $j$, where one takes an arbitrary local element $s$.
For $i>j$, entry $(i,j)$ of the product equals the corresponding
entry of $M(s)$ times
\[
             \prod_{r<j}(d_i-d_r)\prod_{r>j}(d_j-d_r).
\]
This unit multiple vanishes, so all such entries vanish in the same
basis. The matrices remain lower triangular, hence are diagonal
there with the prescribed entries. This recovers the three local
equations from the ordered ordinary condition in any $Q$.
% END G67T2

% BEGIN G67T3
\paragraph{G67T3. Realizing the lower matrices as points of the original entry algebra.}
For a nonzero commutative $k$-algebra $Q$, every pair satisfying
$F.\mathrm{Low}(Q,y,a)$ determines a unique map $A\to Q$
over the residue map on $O$, killing $\mathfrak m_B A$ and
giving exactly $M_{y,a}$ on the universal matrix. The construction
also works for the polynomial targets below, where evaluating an
arbitrary formal series would be unjustified.

Give the adjacent entries of $M_{y,a}$ a factor $t$ and the corner
entry a factor $t^2$. The group law, polarization and normalization
still hold. Locally the parameters become
$(t\alpha,t\beta,t^2\gamma)$, so the ordinary products still vanish.
Let $C_{\rm pol}\subset Q[t]$ be the graded $k$-algebra generated by all
these scaled lower entries. The locally constant functions have
finite image by compactness of $G_0$, so $C_{\rm pol}$ is finitely generated.
Its ideal $I$ of positive degrees has quotient $k$, and
$I^n\subset t^nQ[t]$. Thus $C_{\rm pol}$ injects into its $I$-adic
completion $C$. This is a complete Noetherian local coefficient
algebra with residue field $k$ and $O$-action by reduction.
All normalized matrix equations hold in $C_{\rm pol}$ since it embeds
in $Q[t]$. In particular the local characters are already the
constants $p_s,1,p_s^{-1}$. Over $C$ we therefore have a continuous
lift of the original diagonal residual representation, as the
entries are locally constant and all lower entries belong to $IC$.
Universality supplies $f:R\to C$.

Each generator of the dense subalgebra of $B$ has constant image
in $k$. Indeed every nonvanishing off-diagonal entry has negative
weight, so a weight-zero monomial involving one needs a positive
weight factor and that upper entry vanishes. The diagonal entries
and character generators already belong to $k$. By adic continuity
and closedness of the constant field in $C$ this gives $f(B)\subset k$.
Reduction then identifies $f|_B$ with the prescribed residue.
As $A$ is generated by $B$ and matrix entries, $f(A)\subset C_{\rm pol}$.
Evaluating these polynomials at $t=1$ gives the desired map.
Conversely, RTC10 makes every such fiber map agree on $B$ and
determines it from its entries. Thus this is precisely the map
called $\mathrm{At}$ in the source. Its uniqueness also makes it
compatible with any $k$-algebra map on the target.

In particular every $\mathfrak p\in\mathrm{BB}$ is obtained this
way with $K=\kappa(\mathfrak p)$. The residue map factors through
$k$ on $O$ and kills $\mathfrak m_B A$. Some negative-weight
entry survives by the definition of $\mathrm{BB}$. If an upper
entry survived as well, positive powers of the two entries could
be chosen with weights summing to zero. Their product in $B$
has zero residue, so belongs to $\mathfrak m_B$. Its image in $K$
cannot vanish, a contradiction. Hence every upper entry is zero,
and RTC10 gives the remaining diagonal characters and local constancy.
G67T2 recovers $y,a$ with $F.\mathrm{Low}(K,y,a)$.
% END G67T3

% BEGIN G67T4
\paragraph{G67T4. Nonemptiness, preconnectedness and density off the exceptional ideal.}
By RTC12 choose $0\ne b\in F.\mathrm{strS}(k)$. Its defining
identities are the equations of G67T2 with $y=0$, including
$b_s=(p_s-p_s^{-1})\gamma_v$ locally. Applying G67T3 to $(0,b)$
gives a fixed $\mathfrak p_0\in\mathrm{BB}$, since some corner
entry survives. For any lower pair $(y,a)$, scalar $z$ and strict
corner function $b'$, the pair $(z y,z^2 a+b')$ again satisfies
the lower equations. In the product identity the mixed term scales
by $z^2$. The local parameters become
$(z\alpha,z\beta,z^2\gamma+\gamma'_v)$, and the marked and
polarization equations are preserved.

Let $\mathfrak p$ be any point of $\mathrm{BB}$, and take its
$K$-valued coordinates $(y,a)$ from G67T3. Write $X$ for an
indeterminate and base change the functions along $k\subset K$.
The pair
\[
              (X y,\ X^2a+(1-X)b)
\]
over $K[X]$ thus gives a map $A\to K[X]$. Its prime kernel
$\eta$ contains $\mathfrak m_B A$. Evaluation at $1$ recovers
the map for $\mathfrak p$, and evaluation at $0$ recovers the
map for $\mathfrak p_0$ followed by $k\hookrightarrow K$, by
uniqueness in G67T3. Thus $\eta\subseteq\mathfrak p,\mathfrak p_0$.
In particular $\eta\in\mathrm{BB}$ since
$\mathfrak t_+\not\subseteq\mathfrak p$. The three-point subspace
$\{\eta,\mathfrak p,\mathfrak p_0\}$ has the dense point $\eta$
and hence is connected. The union over all $\mathfrak p$ is
preconnected because these subspaces share $\mathfrak p_0$.

Finally suppose $I_E\subseteq\mathfrak p$. The two adjacent lower
entries have odd weight, so $y=0$ and the corner $a$ is strict.
Use the nonzero adjacent point $(y_*,a_*)$ of G67T1 and take the
pair $(X y_*,X^2a_*+a)$ over $K[X]$. Its prime kernel $\eta'$
specializes at zero to $\mathfrak p$. Hence $\eta'\in\mathrm{BB}$
by the preceding ideal argument. Some adjacent universal entry
maps to the nonzero polynomial $Xy_*(g)$, so
$I_E\not\subseteq\eta'$. Every point of $\mathrm{BB}\cap V(I_E)$
therefore lies in the closure of $\mathrm{BB}\setminus V(I_E)$.
This gives the three conclusions of the original
$\mathrm{preconnected\_rank\_three}$ for every supplied rank-three
space on the frame.
% END G67T4
